% A_appendix.tex
\appendix
\section*{Appendix}

\subsection{Full Proofs}

This appendix contains the full proofs of the theoretical results stated in the main text.

\begin{theorem}[LQRL Reduces Q-Penalty for Improving Skills]
Let a skill $s$ fail a task ($r_{\text{task}} = 0$) but improve its content such that $r_{\text{learning}} = q > 0$. Under the LQRL update with learning rate $\alpha$ and layering weight $\beta$, the Q-value change is:
\begin{equation}
    \Delta Q = \alpha \cdot \beta \cdot q
\end{equation}
Under standard Q-learning ($\beta = 0$), the Q-value change is $\Delta Q_{\text{std}} = -\alpha \cdot Q_{\text{old}} \leq 0$.
\end{theorem}

\begin{proof}
From \cref{eq:lqrl} with $r_{\text{task}} = 0$ and $r_{\text{learning}} = q$:
\begin{align}
    Q_{\text{new}} &= Q_{\text{old}} + \alpha \cdot [(1-\beta) \cdot 0 + \beta \cdot q] \\
    &= Q_{\text{old}} + \alpha \cdot \beta \cdot q
\end{align}
Thus $\Delta Q = \alpha \beta q > 0$ when $q > 0$ and $\beta > 0$.

For standard Q-learning with $r_{\text{task}} = 0$:
\begin{align}
    Q_{\text{new}}^{\text{std}} &= Q_{\text{old}} + \alpha \cdot (0 - Q_{\text{old}}) \\
    &= (1 - \alpha) \cdot Q_{\text{old}}
\end{align}
Thus $\Delta Q_{\text{std}} = -\alpha \cdot Q_{\text{old}} \leq 0$ since $Q_{\text{old}} \geq 0$ and $\alpha > 0$.
\end{proof}

\begin{prop}[Optimal $\beta$ Bound]
For a failed task with content improvement quality $q$, the Q-value under LQRL exceeds that under standard Q-learning when $\beta > \frac{Q_{\text{old}}}{q}$.
\end{prop}

\begin{proof}
LQRL Q-value after failure with improvement: $Q_{\text{LQRL}} = Q_{\text{old}} + \alpha \beta q$.
Standard Q-learning: $Q_{\text{std}} = (1-\alpha) Q_{\text{old}}$.

We want $Q_{\text{LQRL}} > Q_{\text{std}}$:
\begin{align}
    Q_{\text{old}} + \alpha \beta q &> (1-\alpha) Q_{\text{old}} \\
    \alpha \beta q &> Q_{\text{old}} - \alpha Q_{\text{old}} \\
    \beta &> \frac{Q_{\text{old}}(1 - \alpha)}{\alpha q}
\end{align}
For small $\alpha$, this approximates $\beta > \frac{Q_{\text{old}}}{q}$.
\end{proof}

\subsection{Implementation Details}

The LQRL implementation extends the BatchSkillIntegrator in the qskill framework. When a task fails, the system:
\begin{enumerate}
    \item Constructs a failure scenario from the error trace
    \item Evaluates the failure scenario quality using an LLM judge
    \item Updates the skill's failure scenarios with quality control (cap at 10, merge similar)
    \item Computes $r_{\text{learning}}$ based on whether content was added/merged
    \item Applies the LQRL update rule with the computed rewards
\end{enumerate}

The LLM judge prompt asks: ``Did the skill content improve, degrade, or stay the same after this task?'' with criteria for relevance, actionability, non-redundancy, and correctness.
